% Suppress volatile PDF timestamps and trailer IDs so identical sources and
% toolchains produce byte-identical output.
\pdfinfoomitdate=1
\pdftrailerid{}
\documentclass[runningheads]{llncs}

\usepackage[T1]{fontenc}
\usepackage{booktabs}
\usepackage{tabularx}
\usepackage{listings}
\usepackage{xcolor}
\usepackage{hyperref}

\hypersetup{
  colorlinks=true,
  linkcolor=blue!55!black,
  citecolor=blue!55!black,
  urlcolor=blue!55!black,
  pdftitle={Bounded Model Checking, Mechanized Invariants, and Conformance Testing for a Price-Time Matching Engine},
  pdfauthor={Ara Aslyan},
  pdfsubject={A multi-artifact formal-verification case study},
  pdfkeywords={TLA+, Lean 4, model checking, theorem proving, conformance testing, matching engine}
}

\lstset{basicstyle=\ttfamily\scriptsize,breaklines=true,frame=single,
  columns=fullflexible,aboveskip=4pt,belowskip=4pt}
\newcolumntype{Y}{>{\raggedright\arraybackslash}X}
\emergencystretch=1em

\title{Bounded Model Checking, Mechanized Invariants, and Conformance Testing for a Price-Time Matching Engine}
\titlerunning{Verification of a Price-Time Matching Engine}

\author{Ara Aslyan}
\authorrunning{A. Aslyan}
\institute{Strike Technologies, New York, NY, USA\\
\email{aaslyan@striketechnologies.com}}

\begin{document}
\maketitle

\begin{abstract}
We report a case study applying three assurance techniques to related
artifacts for a feature-rich price-time matching engine. A TLA+ model is
exhaustively checked within finite configurations; a Lean~4 executable
reference carries machine-checked invariant-preservation proofs; and traces
from the model are replayed against a C++ implementation on projected
observables. The methods exposed different defects: TLC produced a
three-order FIFO counterexample; operational modeling revealed an ambiguity
between self-trade decrement and iceberg reload; proof development exposed an
insufficient constant bound in the inner matching loop and an MTL/MinQty
transcription error; and conformance testing found a C++ disposal defect.
Four completed TLC configurations explored 55,439,125 distinct states in
aggregate. Lean proves preservation of all modeled book-state invariants for
arbitrary finite books satisfying explicit hypotheses, plus unconditional
post-only and self-trade-prevention guarantees on emitted trades. A named
corollary lifts the result to finite well-formed order sequences from the
empty book. These results do not constitute an end-to-end refinement proof of
the C++ engine. Raw model-checking logs, theorem sources, tests, and
reproduction scripts accompany the artifact.

\keywords{Model checking \and Lean 4 \and TLA+ \and conformance testing
\and matching engine.}
\end{abstract}

\section{Introduction}

A continuous limit-order-book matching engine maintains bids and asks and
selects trades according to price and arrival priority. Even a compact engine
must compose order types, time-in-force rules, partial fills, iceberg reload,
post-only handling, self-trade prevention (STP), minimum quantity, stop
triggering, and cancellation. Local rules are straightforward; their
interactions are not. A triggered stop re-enters the ordinary pipeline, an
STP decrement changes quantities without a trade, and a market-to-limit (MTL)
order changes type after its first execution.

This paper evaluates a layered workflow on three separately written
artifacts: an executable TLA+ state machine, an executable Lean reference,
and a C++ implementation. Each layer answers a different question:

\begin{itemize}
\item TLC checks every reachable state within explicitly bounded TLA+
configurations and returns traces for violations.
\item Lean checks proofs for arbitrary finite data satisfying stated
hypotheses, beyond TLC's size bounds.
\item Trace replay tests agreement between projected TLA+ observations and
the C++ engine on the traces actually exercised.
\end{itemize}

The contribution is an applied case study rather than a new matching
algorithm or proof technique. Its value lies in the defect witnesses,
assurance boundaries, and reproducible cross-artifact checks. The main
findings are:

\begin{enumerate}
\item TLC found a FIFO violation caused by a triggered stop retaining its
original timestamp. The counterexample has been reproduced under a valid
three-order configuration and archived with its raw log.
\item Modeling exposed a gap in the interaction of STP \textsc{decrement}
and iceberg visibility: decrement can exhaust the visible slice without
executing the fill path that reloads it.
\item Lean proof development showed that a constant bound of 100 was
insufficient for the inner matching loop and exposed an MTL/MinQty routing
error in the executable transcriptions.
\item TLA+-trace replay found a C++ path on which a triggered stop-market
remainder could rest on the book.
\end{enumerate}

The artifacts deliberately remain separate. Agreement is tested, not proved,
and correspondence between the prose requirements and each formalization
remains a validation obligation.

\section{Engine and Properties}

An order contains an identifier, side, type, quantity, optional limit and stop
prices, time in force, optional display and minimum quantities, a post-only
flag, and optional STP group and policy. The model supports limit, market,
MTL, stop-limit, and stop-market orders; GTC, IOC, FOK, and DAY; iceberg
display; and four STP policies. Prices are integers drawn from a finite set in
TLC and natural numbers in Lean.

The central state consists of price-indexed bid and ask queues, dormant stops,
last trade price, clock, and next identifier. Price levels are ordered best
first and queues are ordered by timestamp. Processing performs validation,
post-only and FOK/MinQty checks, stop placement or matching, disposal of a
remainder, and a bounded stop-trigger cascade.

Table~\ref{tab:inv} records the canonical property inventory. ``Checked'' in
the TLC column means explicitly listed as an invariant in each checked
configuration; it does not generalize beyond the configuration. The Lean
column separates the capstone from properties proved at another level.

\begin{table}[t]
\caption{Property coverage. INV-10 is outside both formal models.}
\label{tab:inv}
\centering
\scriptsize
\begin{tabularx}{\textwidth}{@{}lYll@{}}
\toprule
ID & Property & TLC & Lean \\
\midrule
1 & No empty price levels & checked & capstone \\
2--3 & Bid/ask level ordering & construction & \texttt{AllInv} \\
4 & Book uncrossed & checked & \texttt{AllInv} \\
5 & No zero-quantity resting orders & checked & capstone \\
6 & Resting-order status consistency & checked & capstone \\
7 & FIFO within a price level & checked & capstone \\
8 & No resting market orders & checked & capstone \\
9 & Execution at passive price & construction & matching-loop theorem \\
10 & Causal event ordering & not modeled & not modeled \\
11 & No trade from post-only aggressor & checked & unconditional \\
12 & No trade within an STP group & checked & unconditional \\
13 & No resting MTL orders & checked & capstone \\
14 & No resting MinQty marker & checked & capstone \\
\bottomrule
\end{tabularx}
\end{table}

\section{Bounded Model Checking in TLA+}

The TLA+ artifact represents books as sequences of price-level records and
orders as records. \texttt{SubmitOrder}, cancellation, amendment, time
advance, matching, disposal, and triggered-stop processing form the next-state
relation. Order shapes are generated once and filtered by the same
well-formedness predicate used by submission. A differential executable check
confirms that Lean and TLA+ accept the same 3,830 bounded order shapes.

TLC's bounds cover submission count, quantity, prices, and logical clock.
Terminal states caused by the clock ceiling are expected; runs therefore
disable deadlock reporting while retaining invariant checks. Table~\ref{tab:tlc}
contains freshly reproduced runs with one worker. The sum of distinct counts
measures aggregate exploration workload, not a deduplicated union across
configurations.

\begin{table}[t]
\caption{Archived TLC results (TLC 2.19, OpenJDK 21, one worker).}
\label{tab:tlc}
\centering
\scriptsize
\begin{tabular}{@{}lrrrrrr@{}}
\toprule
Config. & Orders & Qty & Prices & Generated & Distinct & Time \\
\midrule
Tiny & 2 & 1 & 2 & 2,463,194 & 1,427,827 & 0:32 \\
Small & 2 & 2 & 2 & 16,920,722 & 9,230,323 & 4:12 \\
Medium & 2 & 2 & 3 & 49,656,516 & 29,622,636 & 14:24 \\
Amend & 2 & 2 & 2 & 37,321,458 & 15,158,339 & 25:45 \\
\midrule
Aggregate & & & & 106,361,890 & 55,439,125 & \\
\bottomrule
\end{tabular}
\end{table}

\subsection{FIFO counterexample}

Triggered stops originally retained the timestamp assigned when they became
dormant. Consider three submissions: (1) a buy limit of quantity one at price
one; (2) a sell stop-limit of quantity one with stop and limit price one; and
(3) a sell limit of quantity two at price one. Order~3 fills order~1, setting
the last trade price and triggering order~2. The remaining unit of order~3
rests first with timestamp~3. The triggered order is appended afterward but
retains timestamp~2, yielding the ask queue $[3,2]$ and violating strict FIFO.

With only the timestamp-refresh line reverted, TLC reproduced this violation
under \texttt{MAX\_ORDERS=3}, \texttt{MAX\_QTY=2}, and two prices. It found
9,160,154 generated and 4,980,329 distinct states before the violation in
2m41s on the evaluation host. This is a three-order counterexample; no
minimality claim is made. The repair assigns the current clock when the stop
converts. The raw trace records the seed, fingerprint seed, tool versions,
host, and full state sequence.

\subsection{Iceberg/STP specification gap}

For STP \textsc{decrement}, no trade is emitted; quantities on both orders are
reduced. If the resting order is an iceberg and its visible quantity reaches
zero while hidden quantity remains, a fill-only reload rule never runs. The
result is an unmatchable visible slice of zero. This issue was identified
during operational modeling, not by a pre-fix TLC invariant trace. The model
reloads such an iceberg, refreshes its timestamp, and moves it to the queue
tail. We retain this distinction because a modeled scenario is not equivalent
to an archived model-checker counterexample.

\section{Mechanized Preservation in Lean 4}

The Lean project has no external library dependencies. Its executable
definitions mirror the processing phases while favoring proof-oriented data
structures. \texttt{AllInv} combines uncrossedness with descending bid and
ascending ask levels. \texttt{BookInvariant} adds no empty levels, no ghosts,
status consistency, FIFO, and exclusion of resting market, MTL, and MinQty
orders.

The single-step capstone is:

\begin{lstlisting}
theorem process_preserves_BookInvariant
    (b : BookState) (o : Order)
    (hall : AllInv b) (hpok : OrderProcOk o)
    (hsnp : StopsNoPostOnly b) (hb : BookOk b)
    (hstops : StopsWF b) (hok : OrderRestOk o) :
    BookInvariant (process b o).book
\end{lstlisting}

The hypotheses expose rather than hide the assurance boundary. \texttt{BookOk}
and \texttt{StopsWF} capture structural well-formedness; the order predicates
follow from ordinary order well-formedness. Auxiliary theorems show that the
book predicates, \texttt{AllInv}, and \texttt{StopsNoPostOnly} hold initially
and are preserved. The named execution-level corollary packages the induction:

\begin{lstlisting}
theorem process_all_preserves_BookInvariant (orders : List Order)
    (hwf : forall o in orders, OrderProcOk o /\ OrderRestOk o) :
    BookInvariant
      (orders.foldl processBook BookState.empty)
\end{lstlisting}

Thus the theorem covers every finite submitted-order sequence satisfying the
per-order predicates, not unconstrained values. Separate unconditional
theorems establish post-only and STP safety for the emitted trade list. Passive
pricing is proved at the matching-loop level but is not part of the capstone;
event ordering is not represented.

\subsection{Fuel and proof structure}

The original inner matching loop used the constant 100. This made a universal
preservation statement false for books with more than 100 matchable opposing
orders: fuel could expire while a crossing book remained. The repair computes
fuel from the opposing side and proves that every recursive step decreases the
relevant order count or incoming quantity. This sufficiency result applies to
the inner \texttt{doMatch} loop. The outer order/stop-cascade worker likewise
uses no fixed limit: \texttt{computeProcessFuel} derives its budget from both
book sides, all dormant stops, and the incoming order. Thus neither executable
processing layer contains an arbitrary constant cutoff. The safety theorems
are additionally parametric in the internal worker budget.

The constructive development proves preservation branch by branch. A shorter
structural file gives a complementary uncrossedness proof using the matching
termination trichotomy, but imports the constructive fuel result and is not an
independent assurance argument. The complete reachable-state theorem is in a
fourth file. All are imported by the library root, so a default build
elaborates them.

The verification gate rejects \texttt{sorry}, prints the axiom dependencies of
ten top-level theorems, and runs the executable tests. The inspected results
depend only on \texttt{propext}, \texttt{Classical.choice}, and
\texttt{Quot.sound} as applicable; tactics such as \texttt{omega} produce
terms checked by the Lean kernel.

\section{Trace-Based C++ Conformance}

The C++ engine is separately implemented with STL containers. Conformance is
tested by generating short TLA+ target-search traces, converting submissions
and expected observations to JSON, replaying the submissions, and comparing
top of book, order and stop counts, last trade price, and fills after each
step. These are projected observations, not full-state equivalence.

A fresh fast sweep ran 117 target searches: 13 seeded scenarios by nine
structural targets. Ninety-four convertible traces were replayed, comprising
109 observation steps, 384 submitted orders, and 40 fills; all replayed traces
passed with no divergence. Some target searches terminate without the named
target but still yield a syntactically convertible TLC behavior, which is why
the number of converted traces can exceed the 57 runs that reached their
target. Twenty-three runs yielded no convertible behavior. The safety phase
used a deliberately short 120-second cap: four scenarios completed and nine
timed out without a violation in the explored prefix. These timeouts are not
reported as passes.

The fast sweep is evaluation evidence, while the verification gate contains a
smaller deterministic regression set (10 traces, 21 steps) plus 65 directed
C++ tests and a 2,000-step shadow differential. The historical C++ defect was
a disposal path allowing a persistent triggered stop-market remainder to
rest. Trace replay localized the mismatch; after repair, market remainders are
canceled rather than inserted.

Conformance cannot establish correctness for untested paths, and two
implementations can agree on the same error. A specification-rooted oracle is
useful precisely when it differs from the implementation's design lineage,
but it still supplies testing rather than refinement.

\section{Cross-Layer Lessons and Limitations}

The layers found different faults because their quantifiers differ. TLC
explores complete bounded executions and produces witnesses. Lean quantifies
over arbitrary finite structures satisfying hypotheses, exposing bounds that
small configurations never reach. Conformance testing addresses transcription
and implementation gaps but only for sampled traces. No one result subsumes
the others.

The artifact has several explicit limits. TLC prices, quantities, submissions,
and time are finite; only one non-null STP group identifier is used; GTD timing
is simplified; and event ordering is absent. The Lean theorem is about the
Lean executable and requires well-formed inputs and state. There is no
Lean-to-C++ or TLA+-to-C++ refinement theorem.
The prose-to-formal correspondence is reviewed and differentially tested but
not proved. The MTL/MinQty episode itself demonstrates why this distinction
matters: a passing precheck could take an erroneous disposal route in both
formal executable transcriptions until the full invariant proof forced that
case to be resolved.

Threats to validity include the small TLC domains, model and implementation
sharing terminology, incomplete long-running conformance safety searches, and
hardware-dependent timing. We therefore archive raw counts and distinguish a
completed exploration, a violation-terminated run, a timeout, and a replay
pass. Aggregate distinct-state counts are workload totals rather than unique
states across models.

\section{Related Work}

Industrial TLA+ case studies demonstrate the value of executable
specifications and counterexample traces~\cite{Newcombe2015,Newcombe2014}; our
focus is their composition with a mechanized executable and implementation
replay. Model-based testing provides the broader basis for deriving tests from
transition systems~\cite{Tretmans2008,Utting2012}.

Mechanized financial-market work includes verified trades and double-auction
formalizations~\cite{Sarswat2020,Natarajan2021,Garg2025}, while automated
reasoning has been applied directly to financial
algorithms~\cite{Passmore2017,Passmore2021}.
Our contribution differs in emphasizing feature interaction across three
non-refined artifacts rather than a new auction theorem. The Lean termination
argument follows established proof-assistant treatments of recursive
definitions and well-founded measures~\cite{Krauss2009,Bove2016}.

Differential testing detects behavioral disagreement but cannot identify the
correct side or errors shared by both implementations~\cite{McKeeman1998}.
Here the TLA+ projection supplies a third oracle, while the Lean proof supplies
unbounded invariant reasoning for a related executable.

\section{Reproducibility and Conclusion}

The artifact provides a fast gate and a deep gate. The fast gate builds Lean,
checks proof closure and axiom footprints, runs Lean and C++ tests, TLC smoke
and invariant-coverage checks, conformance replay, shadow differential, and
the Lean/TLA+ well-formedness differential. The deep gate reruns the four
completed TLC configurations; a bounded three-order attempt is reported as a
skip unless it completes. Raw logs include commands, hashes, versions, seeds,
queue sizes, and host details. CI runs the fast gate and paper build on changes
and schedules the deep gate.

This case study supports a deliberately delimited conclusion: bounded TLA+
exploration, Lean invariant preservation, and model-based C++ conformance
testing expose complementary defect classes. Lean proves the modeled
book-state properties along finite well-formed executions of its reference;
TLC exhausts only stated finite configurations; and C++ agreement is tested on
replayed projections. Together they provide stronger evidence than any layer
alone, but not an end-to-end proof of the deployed implementation.

\subsubsection{Data availability.}
Source, proof artifacts, configurations, raw logs, and scripts are available
in the \href{https://github.com/formally-verified-exchange/verified-matching-engine}{project repository}
under the MIT license. The repository records SHA-256 hashes for archived
model-checking inputs. A versioned archival DOI should replace this repository
URL in the camera-ready version once the release is deposited.

\subsubsection{AI assistance.}
The formal artifacts and manuscript were developed interactively with a large
language model used as a coding and drafting assistant. Proof terms are
accepted by the Lean kernel, and empirical claims are tied to archived tool
outputs; responsibility for the claims and artifact remains with the author.

\begin{thebibliography}{20}

\bibitem{Bove2016}
Bove, A., Krauss, A., Sozeau, M.: Partiality and recursion in interactive
theorem provers---an overview. Math. Struct. Comput. Sci. 26(1), 38--88 (2016)

\bibitem{Garg2025}
Garg, M., Natarajan, R., Sarswat, S., Singh, A.K.: Double auctions:
formalization and automated checkers. J. Autom. Reason. 69, 17 (2025)

\bibitem{Krauss2009}
Krauss, A.: Automating Recursive Definitions and Termination Proofs in
Higher-Order Logic. Ph.D. thesis, Technische Universit\"at M\"unchen (2009)

\bibitem{McKeeman1998}
McKeeman, W.M.: Differential testing for software. Digital Technical Journal
10(1), 100--107 (1998)

\bibitem{Natarajan2021}
Natarajan, R., Sarswat, S., Singh, A.K.: Verified double sided auctions for
financial markets. In: ITP 2021. LIPIcs, vol. 193, pp. 28:1--28:18 (2021)

\bibitem{Newcombe2014}
Newcombe, C.: Why Amazon chose TLA+. In: ABZ 2014. LNCS, vol. 8477,
pp. 25--39. Springer (2014)

\bibitem{Newcombe2015}
Newcombe, C., Rath, T., Zhang, F., Munteanu, B., Brooker, M., Deardeuff, M.:
How Amazon Web Services uses formal methods. Commun. ACM 58(4), 66--73 (2015)

\bibitem{Passmore2017}
Passmore, G.O., Ignatovich, D.: Formal verification of financial algorithms.
In: CADE 2017. LNCS, vol. 10395, pp. 26--41. Springer (2017)

\bibitem{Passmore2021}
Passmore, G.O.: Some lessons learned in the industrialization of formal
methods for financial algorithms. In: FM 2021. LNCS, vol. 13047,
pp. 717--721. Springer (2021)

\bibitem{Sarswat2020}
Sarswat, S., Singh, A.K.: Formally verified trades in financial markets.
In: ICFEM 2020. LNCS, vol. 12531, pp. 217--232. Springer (2020)

\bibitem{Tretmans2008}
Tretmans, J.: Model based testing with labelled transition systems. In:
Formal Methods and Testing. LNCS, vol. 4949, pp. 1--38. Springer (2008)

\bibitem{Utting2012}
Utting, M., Pretschner, A., Legeard, B.: A taxonomy of model-based testing
approaches. Softw. Test. Verif. Reliab. 22(5), 297--312 (2012)

\end{thebibliography}

\end{document}
